\section{特征标是表示的指纹}
\label{sec:04-Discrete-Mathematics/09-Representation-Theory/03}

有人递给你一个 \(S_3\) 的六维表示，六个 \(6\times6\) 矩阵，全都写满了非零元素。问：它由哪些不可约成分拼成，各几份？

按上一节的定义硬做，要在 \(\C^6\) 里搜索不变子空间。可这些矩阵长什么样，取决于对方当初选了哪组基；换一组基，矩阵面目全非，而分解结果一个字都不会变。既然答案与基无关，就不该把力气花在依赖基的对象上。

需要一个在相似变换下岿然不动的量。这样的量并不多，最好用的一个是迹：\(\trace(TAT^{-1})=\trace A\)。

\subsection{一个共轭类上一个数}

\begin{definition}
表示 \(\rho\) 的特征标是函数
\[
\chi_\rho(g)=\trace\rho(g),\qquad g\in G.
\]
\end{definition}

两条性质当场可得。等价的表示有相同的特征标，因为迹不认基。而且 \(\chi_\rho(hgh^{-1})=\chi_\rho(g)\)，所以特征标是类函数——它在每个共轭类上取常值，实际信息量只有``共轭类个数''这么多。\(\chi_\rho(e)=\trace I=\dim V\)，维数是白送的第一个数。

一个 \(n\) 维表示原本要 \(|G|\cdot n^2\) 个数才写得完，压缩成共轭类个数那么多个数之后，居然没有丢掉要紧的东西：复数域上，两个表示等价当且仅当特征标相同。这一步压缩是本节全部计算便利的来源。

\subsection{正交关系把分解变成内积}

在类函数空间上放一个内积
\[
\ip{\varphi}{\psi}=\frac1{|G|}\sum_{g\in G}\overline{\varphi(g)}\,\psi(g).
\]

\begin{theorem}[特征标正交关系]
设 \(\chi_1,\ldots,\chi_r\) 是有限群 \(G\) 的全部不可约复特征标，则 \(\ip{\chi_i}{\chi_j}=\delta_{ij}\)。它们构成类函数空间的一组标准正交基。
\end{theorem}

证明沿的是上一节末尾那条线：群代数分解为矩阵代数的直和，Schur 引理迫使不同块之间的交换映射为零，取迹后就得到正交性。求和的技术细节可以跳过，因为结论本身已经把用法交代完了。

既然每个表示都是不可约表示的直和，特征标就是相应特征标的和，\(\chi_\rho=\sum_i m_i\chi_i\)。正交基一来，系数用一次内积读出：
\[
m_i=\ip{\chi_i}{\chi_\rho}.
\]
这与在正交基下求坐标是同一个动作，只不过向量换成了群上的函数。求出来的 \(m_i\) 必须是非负整数，这本身还是一个检验：算出小数或负数，说明前面某处错了。顺带还得到一个判据——\(\ip{\chi_\rho}{\chi_\rho}=1\) 当且仅当 \(\rho\) 不可约，因为平方和等于 \(1\) 的非负整数组只有一个 \(1\)。

\begin{center}
\begin{tikzpicture}[x=1cm,y=1cm,font=\footnotesize,>=stealth]
  \draw[black!25,step=1cm] (0,0) grid (3,3);
  \draw[black!65] (0,0) rectangle (3,3);
  \node[black!70] at (0.5,3.35) {\(\{e\}\)};
  \node[black!70] at (1.5,3.35) {对换};
  \node[black!70] at (2.5,3.35) {三轮换};
  \node[left,black!70] at (-0.15,2.5) {平凡};
  \node[left,black!70] at (-0.15,1.5) {符号};
  \node[left,black!70] at (-0.15,0.5) {二维};
  \node at (0.5,2.5) {\(1\)}; \node at (1.5,2.5) {\(1\)}; \node at (2.5,2.5) {\(1\)};
  \node at (0.5,1.5) {\(1\)}; \node at (1.5,1.5) {\(-1\)}; \node at (2.5,1.5) {\(1\)};
  \node at (0.5,0.5) {\(2\)}; \node at (1.5,0.5) {\(0\)}; \node at (2.5,0.5) {\(-1\)};

  \draw[->,black!70] (5.0,0.3) -- (8.6,0.3) node[right,black!80] {\(\chi_1\)};
  \draw[->,black!70] (5.0,0.3) -- (5.0,3.3) node[above,black!80] {\(\chi_2\)};
  \draw[->,very thick,blue!55!black] (5.0,0.3) -- (7.7,2.5);
  \node[above,blue!55!black] at (7.5,2.5) {\(\chi_\rho\)};
  \draw[dashed,black!45] (7.7,2.5) -- (7.7,0.3);
  \draw[dashed,black!45] (7.7,2.5) -- (5.0,2.5);
  \node[below,black!70] at (7.7,0.25) {\(m_1\)};
  \node[left,black!70] at (4.95,2.5) {\(m_2\)};
\end{tikzpicture}

\emph{左：\(S_3\) 的特征标表，每一行是一个不可约表示的指纹，每个共轭类占一列。右：这些指纹在类函数空间中两两正交，任一表示的特征标落在它们张成的格点上，坐标 \(m_i\) 就是一次内积的结果。找不变子空间的搜索问题，被换成了求坐标。}
\end{center}

\subsection{正则表示为什么把每个原子都装了进去}

拿正则表示试一次。\(\C[G]\) 上 \(\rho(g)\) 把基向量 \(h\) 送到 \(gh\)，对角元非零当且仅当 \(gh=h\)，即 \(g=e\)。所以
\[
\chi_{\text{reg}}(e)=|G|,\qquad \chi_{\text{reg}}(g)=0\ (g\ne e).
\]
求内积时只有 \(g=e\) 一项活着，
\[
m_i=\frac1{|G|}\,\overline{\chi_i(e)}\,\chi_{\text{reg}}(e)=\overline{d_i}=d_i.
\]
每个不可约表示都在正则表示里出现，出现次数正好等于自己的维数——顺便把上一节的平方和公式 \(\sum_i d_i^2=|G|\) 又验证了一遍。

对 \(S_3\) 就是 \(6=1+1+2\times2\)。开头那个六维表示如果恰好是正则表示，答案便是：平凡与符号各一份，二维那个两份。

\(D_4\) 的表大一点，五个共轭类五个不可约表示，二十五个数把这个群的表示论装完了。

{\def\LTcaptype{none}
\begin{longtable}{@{}lrrrrr@{}}
\toprule\noalign{}
 & \(e\) & \(r^2\) & \(r,r^3\) & \(s,r^2s\) & \(rs,r^3s\) \\
\midrule\noalign{}
\endhead
\bottomrule\noalign{}
\endlastfoot
类的大小 & 1 & 1 & 2 & 2 & 2 \\
\(A_1\)（平凡） & 1 & 1 & 1 & 1 & 1 \\
\(A_2\) & 1 & 1 & 1 & \(-1\) & \(-1\) \\
\(B_1\) & 1 & 1 & \(-1\) & 1 & \(-1\) \\
\(B_2\) & 1 & 1 & \(-1\) & \(-1\) & 1 \\
\(E\)（二维） & 2 & \(-2\) & 0 & 0 & 0 \\
\end{longtable}
}

读表时把``类的大小''那一行当作内积的权重：求和跑遍群元素，同一类里的值相同，于是每一列贡献一次乘积再乘上类的大小。核对 \(\ip{\chi_{A_1}}{\chi_{A_2}}=0\) 只需五项相加，几秒钟的事，这里不占篇幅。真正值得注意的是最后一行：\(E\) 在三个类上取零，而 \(\chi(e)=2\) 与 \(\chi(r^2)=-2\) 说明 \(r^2\) 在二维块上作用为 \(-I\)。特征标不只用来做分解，它还能反过来讲出群元素在各块上干了什么。

\subsection{交换群上的表示论就是 Fourier 分析}

\(G\) 交换时每个共轭类只有一个元素，共轭类数等于 \(|G|\)，于是不可约表示有 \(|G|\) 个，平方和公式逼得它们全是一维的。一维表示与自己的特征标是同一个东西，此时``特征标''退化为群到 \(\C^\times\) 的同态。

取 \(G=\Z/n\Z\)。同态由 \(1\) 的去向决定，而它必须映到一个 \(n\) 次单位根，所以
\[
\chi_k(j)=e^{2\pi\mathrm{i}jk/n},\qquad k=0,1,\ldots,n-1.
\]
这些正是离散 Fourier 变换的基函数。正交关系 \(\ip{\chi_k}{\chi_l}=\delta_{kl}\) 就是 DFT 矩阵按 \(1/\sqrt n\) 归一化后成为酉矩阵的那条性质。\hyperref[sec:08-Numerical-Analysis/08-Fourier-and-Spectral-Methods/03]{``DFT 是变换，FFT 是算法''}里那组基不是凑巧挑得漂亮，它是循环群仅有的那些不可约表示。

由此还能解释卷积为什么在频域变成逐点乘法。循环群上的卷积算子与所有平移交换，Schur 引理说这样的算子在每个一维不可约块上只能是标量——那个标量就是 Fourier 系数。同一句话搬到非交换群，块的尺寸变成 \(d_i\)，卷积于是变成逐块的 \(d_i\times d_i\) 矩阵乘法，而不再是逐点相乘。快速 Fourier 变换往非交换有限群上的推广，靠的正是这种分块结构来组织递归。

把这条线拉长一点：\hyperref[ch:068]{``Fourier 与谱方法：把函数换到频率坐标中''}的全部内容，可以读成一个特定群（循环群、圆周群、实数加群）上的表示论。频率不是分析学的偶然发明，它是平移群的不可约表示的编号。

\subsection{有限、复数、完全可约，三条都在承重}

本节所有算式都建立在有限求和之上。换成紧群，比如三维旋转群 \(SO(3)\)，特征标与正交理论仍在，但求和要换成 Haar 测度下的积分，完备性由 Peter--Weyl 定理接管，不可约表示按最高权分类而不再按共轭类计数。有限性一旦松开，就得换一整套分析工具。

底域也在承重。复特征标表不能不加检查地当作实表示的分类结果：一个复不可约表示可能是实型、复型或四元数型，Frobenius--Schur 指标才给出这项判别。晶体学与量子力学里之所以要区分这三型，原因就在这里。

还有一条前面反复出现的边界：域的特征整除 \(|G|\) 时 Maschke 定理失效，表示不再完全可约，``分解成原子''这套说法整体不适用，特征标也不再决定表示。模表示论是另一片水域。

在这三条之内，Part 04 的这段阶梯算是走完了：群给出对称操作，群作用让它落在集合上，表示把它搬进向量空间，不可约表示是拆不动的模式，特征标是识别模式的指纹。往后遇到球谐函数、等变层里那些看似古怪的通道设计、或者卷积定理的推广，回到这条阶梯上找位置，通常比逐个记结论走得远。
